% Λύση του 11ου προβλήματος της ενότητας «Άλυτα Προβλήματα».
\subsection*{Λύση Προβλήματος 11 --- Αρχή ελάχιστου χρόνου του \en{Fermat}}

\begin{alist}
  \item Στο ορθογώνιο τρίγωνο $AA'P$ ισχύει
        \[
          AP=\sqrt{x^2+a^2}.
        \]
        Επειδή $A'B'=d$ και $A'P=x$, έχουμε
        \[
          PB'=d-x.
        \]
        Άρα, από το ορθογώνιο τρίγωνο $PBB'$,
        \[
          PB=\sqrt{(d-x)^2+b^2}.
        \]

        Ο χρόνος κίνησης στην άμμο είναι
        \[
          T_1(x)=\frac{AP}{v_1}
          =\frac{\sqrt{x^2+a^2}}{v_1},
        \]
        ενώ ο χρόνος κολύμβησης είναι
        \[
          T_2(x)=\frac{PB}{v_2}
          =\frac{\sqrt{(d-x)^2+b^2}}{v_2}.
        \]
        Επομένως, ο συνολικός χρόνος είναι
        \[
          \boxed{
          T(x)=\frac{\sqrt{x^2+a^2}}{v_1}
          +\frac{\sqrt{(d-x)^2+b^2}}{v_2}},
          \qquad x\in[0,d].
        \]

  \item Για $x\in(0,d)$ παραγωγίζουμε:
        \[
          T'(x)
          =
          \frac{x}{v_1\sqrt{x^2+a^2}}
          -
          \frac{d-x}{v_2\sqrt{(d-x)^2+b^2}}.
        \]
        Επίσης,
        \[
          T''(x)
          =
          \frac{a^2}{v_1(x^2+a^2)^{3/2}}
          +
          \frac{b^2}{v_2\bigl((d-x)^2+b^2\bigr)^{3/2}}>0.
        \]
        Άρα η $T'$ είναι γνησίως αύξουσα στο $(0,d)$.

        Επιπλέον,
        \[
          T'(0)=-\frac{d}{v_2\sqrt{d^2+b^2}}<0
        \]
        και
        \[
          T'(d)=\frac{d}{v_1\sqrt{d^2+a^2}}>0.
        \]
        Επομένως, υπάρχει μοναδικό $x_0\in(0,d)$ για το οποίο
        $T'(x_0)=0$. Η $T$ είναι γνησίως φθίνουσα πριν από το $x_0$
        και γνησίως αύξουσα μετά από αυτό, οπότε στο $x_0$ παρουσιάζει
        το ολικό της ελάχιστο.

        Από τη συνθήκη $T'(x_0)=0$ προκύπτει
        \[
          \frac{x_0}{v_1\sqrt{x_0^2+a^2}}
          =
          \frac{d-x_0}
          {v_2\sqrt{(d-x_0)^2+b^2}}.
        \]
        Οι γωνίες $\theta_1$ και $\theta_2$ μετρώνται ως προς την
        κάθετη στην ακτογραμμή. Από τα αντίστοιχα ορθογώνια τρίγωνα,
        \[
          \hm\theta_1
          =\frac{x_0}{\sqrt{x_0^2+a^2}}
        \]
        και
        \[
          \hm\theta_2
          =\frac{d-x_0}{\sqrt{(d-x_0)^2+b^2}}.
        \]
        Άρα η συνθήκη ελαχίστου γράφεται
        \[
          \frac{\hm\theta_1}{v_1}
          =
          \frac{\hm\theta_2}{v_2},
        \]
        δηλαδή
        \[
          \boxed{
          \frac{\hm\theta_1}{\hm\theta_2}
          =\frac{v_1}{v_2}}.
        \]
        Επομένως, η διαδρομή ελάχιστου χρόνου ικανοποιεί τον νόμο του
        \en{Snell}.
\end{alist}
